
Evidence-centric robustness concerns the quality, integrity, and sufficiency of the evidence that enters the context window. In classical information retrieval, irrelevant results are often treated as a ranking problem. In RAG, however, the retrieved set becomes part of the model input; a single misleading passage can change the generated answer, confidence, and citations. Evidence-centric work therefore asks three questions: what kinds of bad evidence appear, how do they affect generation, and which pre-generation controls can reduce their harm?

\begin{table*}[t]
\centering
\small
\begin{tabularx}{\linewidth}{p{0.18\linewidth}p{0.18\linewidth}p{0.25\linewidth}X}
\toprule
\textbf{Theme} & \textbf{Representative work} & \textbf{Core idea} & \textbf{What it reveals for robust RAG} \\
\midrule
Irrelevant context & \citetalias{irrelevantContext2023} & Analyze and mitigate irrelevant retrieved context with filtering and robust training & Even answerable questions degrade when irrelevant passages enter the context; robustness is not equivalent to retrieval recall. \\
Retrieval-noise training & \citetalias{raat2024} & Categorize retrieval noise and train with adaptive adversarial objectives & Noise types should be separated: irrelevant, relevant-but-insufficient, and counterfactual evidence have different effects. \\
Generator note-taking & \citetalias{chainOfNote2023} & Generate per-passage notes before answering & Robust generation benefits from intermediate relevance judgments rather than direct answer generation. \\
Fine-grained noise benchmark & \citetalias{rrag027} & Configurable benchmark for distracting, low-quality, inconsequential, and irrelevant noise & ``Noise'' is not monolithic; evaluation should vary type, amount, and placement of bad evidence. \\
Nonlinear distraction & \citetalias{rrag003} & Study disproportionate effect of small amounts of misleading information & A few strong distractors can dominate long-context reasoning, so average noise rate is insufficient. \\
Long-context RAG & \citetalias{rrag039} & Analyze long-input RAG and hard negatives & More retrieved context can hurt unless ordering, filtering, or training adapts to long-context failure modes. \\
Sentence-level filtering & \citetalias{rrag044} & Filter retrieved context before generation & Filtering can be more fine-grained than document reranking and can remove local distractors inside otherwise useful passages. \\
Retrieval evaluation & \citetalias{rrag058} & Utility-aware retrieval metric for LLMs & RAG retrieval should be evaluated by downstream utility, including negative utility from distracting passages. \\
\bottomrule
\end{tabularx}
\caption{Representative evidence-centric work on noisy retrieval and distraction.}
\label{tab:evidence-methods}
\end{table*}


\paragraph{Retrieval Noise: Beyond Relevant vs. Irrelevant}

The earliest robust-RAG studies show that retrieval errors do not merely reduce recall; they alter the generator's behavior. \citet{irrelevantContext2023} demonstrate that retrieval-augmented language models can be brittle to irrelevant context and propose filtering and robust training as mitigation. This work is important because it separates two effects that are often conflated: a retriever may fail to retrieve enough relevant evidence, and a generator may fail to ignore irrelevant evidence once it is retrieved. Robustness therefore depends on both retriever quality and generator selectivity.

Subsequent work expands the typology of noise. \citet{raat2024} distinguishes multiple retrieval-noise types and trains models with adaptive adversarial objectives and explicit noise awareness. The distinction matters: an irrelevant passage that is obviously off-topic may be easy to ignore, while a counterfactual or hard-negative passage may be semantically close to the query and therefore much more dangerous. RGB similarly evaluates noise robustness, negative rejection, information integration, and counterfactual robustness as separate abilities \citep{rgb2023}. Taken together, these works suggest that ``RAG robustness'' cannot be measured by a single noisy-context accuracy number. A strong evaluation should vary the type, quantity, ranking position, and semantic hardness of negative evidence.

Recent benchmarks make this point more explicit. Magic Mushroom introduces fine-grained retrieval noise categories, including distracting and low-quality evidence \citep{rrag027}. URAG formulates uncertainty quantification in RAG and studies the effect of irrelevant and noisy retrieval on conformal prediction \citep{rrag001}. PRGB and other placeholder or perturbation-based benchmarks test whether models remain robust when retrieval context is systematically manipulated \citep{rrag183}. The trajectory of this literature is clear: robust RAG evaluation is moving from randomly injected distractors to controlled stress tests that expose distinct failure mechanisms.

\paragraph{Distraction and Context Utility}

Distraction is stronger than irrelevance. A passage is distracting when it has negative marginal utility: adding it changes the answer, confidence, attribution, or reasoning path in a harmful direction. This explains why more retrieval is not always better. In long-context settings, retrieving more passages can increase the chance that a hard negative appears before or near useful evidence. \citet{rrag039} study long-context RAG and show why ordering and hard negatives become central once the generator can consume many passages. \citet{rrag003}, whose title names the \emph{first drop of ink}, identify a nonlinear effect in which a small amount of misleading information disproportionately degrades long-context reasoning. CUB evaluates context-utilisation techniques across gold, conflicting, and irrelevant contexts \citep{rrag032}, while UDCG redefines retrieval evaluation around downstream utility rather than traditional relevance alone \citep{rrag058}.

This line of work changes the role of the retriever. The retriever is not only a recall engine; it is a context constructor. A retriever optimised for lexical or embedding similarity may return passages that are highly related to the query but harmful to the generator. A robust retriever should consider downstream answer utility, redundancy, conflict, and the interaction between passages. Similarly, a robust RAG system should not treat the top-$k$ list as a fixed input. It should be able to filter, reorder, compress, diversify, or abstain when the context set is unreliable.

\paragraph{Evidence-Side Mitigations}

Evidence-side defenses operate before or during context construction. They include retriever retraining, reranking, filtering, passage compression, and context ordering. Filtering can be document-level or sentence-level. FILCO-style sentence filtering removes irrelevant local content inside otherwise useful passages \citep{rrag044}; this is important because a retrieved document can contain both answer evidence and distracting details. Utility-aware retrieval metrics and rerankers attempt to optimise what helps the generator, not what merely matches the query \citep{rrag058}. Context ordering and long-context tuning address cases where the same evidence set has different effects depending on position \citep{rrag039}.

Training-based approaches expose the generator to noise so that it learns to separate relevant from irrelevant context \citep{raat2024,rrag091}. This can improve robustness but also raises a scalability question: robust training may be less effective or less necessary as base LLMs become stronger, and adversarially generated noise can overfit to a particular retriever or generator. Training-free approaches, such as filtering and note-taking, are easier to add to existing systems but depend on reliable relevance or sufficiency signals. A recurring pattern is that no single pre-generation control is sufficient: retrieval filtering reduces easy noise, reranking handles ordering, compression improves context budget, and downstream calibration remains necessary for residual uncertainty.

\paragraph{Query-Side and Index-Side Sources of Noise}

Evidence noise does not begin at retrieval scoring. Query formulation is itself a brittle interface. Query-level robustness work shows that minor changes in wording can substantially affect retrieval and downstream answers \citep{queryRobustness2025}. This suggests that robust RAG should include query normalisation, query expansion, and query uncertainty estimation. In multi-turn settings, the query may be non-standalone or under-specified, making retrieval quality depend on dialogue state tracking \citep{rrag141}. The retriever may also inherit biases from chunking and indexing: relevant facts can be split across chunks, qualifiers can be separated from claims, and embeddings can overemphasise topical similarity over answer sufficiency.

Index-side design also changes robustness. Smaller chunks improve retrieval granularity but can remove context needed for factual interpretation. Larger chunks preserve context but increase the chance of including distracting material. Hybrid sparse--dense retrieval can reduce some embedding failures but may increase redundancy. Graph-based, multimodal, and agentic/process-level RAG add structure but introduce new poisoning surfaces, as shown by GraphRAG poisoning, graph poisoning, and reasoning-style poisoning work \citep{rrag051,rrag127,rrag175}. Thus, evidence-centric robustness should be studied as a corpus construction problem, not only as a retriever ranking problem.

\paragraph{Design Tradeoffs}

Evidence-side controls have three recurring tradeoffs. First, filtering improves precision but can remove weakly relevant evidence needed for multi-hop reasoning. Second, reranking can demote distractors but may also amplify model-specific biases if the reranker is trained against a narrow generator. Third, context compression saves tokens but can delete qualifiers, source information, or uncertainty markers. Robust RAG systems should therefore expose confidence about filtering decisions and allow later stages to recover from false negatives. A practical design is multi-stage: retrieve broadly, rerank cautiously, filter locally, preserve provenance, and let the generator or verifier request more evidence when sufficiency is low.
